\chapter[Band 33: Halbgruppen mit Nullelement]{Band 33 auf einen Blick}
\noindent\textbf{Halbgruppen mit Nullelement}\par\medskip
Das Nullelement absorbiert auf beiden Seiten. Es ist von einem
neutralen Element zu unterscheiden. Die konkrete Operation bleibt stets Teil
der Struktur.

\begin{BandUebersicht}
\textbf{Beidseitige Nullabsorption}
\UebFormel{\begin{aligned}
&\SemigroupWithZeroAlg A\star0\ \leftrightarrow\\
&\quad\SemigroupAlg A\star\land0\in A\land\\
&\quad\forall a\in A\
 (0\star a=0\land a\star0=0).
\end{aligned}}
\UebQuellen{\FormulaRefAuto{SemigroupWithZeroDef}[def]}
&
\textbf{Eindeutigkeit der Null}
\UebFormel{\begin{aligned}
&\SemigroupWithZeroAlg A\star0,\ z\in A,\\
&\forall a\in A\
 (z\star a=z\land a\star z=z)\ \vdash\\
&\qquad z=0.
\end{aligned}}
\UebQuellen{\FormulaRefAuto{SemigroupZeroUnique}[thm]} \\
\addlinespace[.8ex]

\textbf{Konstante Operation}
Für \(0\in A\) werde die Operation festgelegt durch\UebFormel{\begin{aligned}
&a,b\in A\ \vdash\ a\star b=0.
\end{aligned}}
\UebQuellen{\FormulaRefAuto{SemigroupWithZeroDef}[def]}
&
\textbf{Konstante und natürliche Beispiele}
\UebFormel{\begin{aligned}
&0\in A,\ \forall a,b\in A\ (a\star b=0)\
 \vdash\\
&\qquad\SemigroupWithZeroAlg A\star0,\\
&\SemigroupWithZeroAlg{\mathbb N}{\cdot}{0}.
\end{aligned}}
\UebQuellen{\FormulaRefAuto{ConstantSemigroupWithZero}[thm];
\FormulaRefAuto{NaturalMultiplicationSemigroupWithZero}[thm]} \\
\addlinespace[.8ex]

\textbf{Annihilatoren}
Für \(\SemigroupWithZeroAlg A\star0\) und \(X\subseteq A\)
schreiben wir kurz\UebFormel{\begin{aligned}
&L(X)=\{a\in A\mid\\
&\hspace{4em}\forall x\in X,\ a\star x=0\},\\
&R(X)=\{a\in A\mid\\
&\hspace{4em}\forall x\in X,\ x\star a=0\}.
\end{aligned}}
\UebQuellen{\FormulaRefAuto{SemigroupLeftAnnihilatorDef}[def];
\FormulaRefAuto{SemigroupRightAnnihilatorDef}[def]}
&
\textbf{Zweiseitige Bedingung}
\UebFormel{\begin{aligned}
&\operatorname{Ann}(X)=L(X)\cap R(X),\\
&a\in\operatorname{Ann}(X)\ \leftrightarrow\\
&\quad a\in A\land\forall x\in X\\
&\qquad(a\star x=0\land x\star a=0).
\end{aligned}}Dies entfaltet die drei Annihilatordefinitionen.
\UebQuellen{\FormulaRefAuto{SemigroupAnnihilatorDef}[def];
\FormulaRefAuto{SemigroupLeftAnnihilatorDef}[def];
\FormulaRefAuto{SemigroupRightAnnihilatorDef}[def]} \\
\addlinespace[.8ex]

\textbf{Homomorphismen und Isomorphismen}
Für zwei Strukturen derselben Klasse bedeutet ein Homomorphismus
\(H(f)\):\UebFormel{\begin{aligned}
&f:A\to B,\\
&x,y\in A\ \vdash\\
&\quad f(x\star y)=f(x)\diamond f(y).
\end{aligned}}Zusätzlich gilt \(f(0_A)=0_B\).Ein Isomorphismus \(I(f)\) ist ein bijektiver Homomorphismus:\UebFormel{\begin{aligned}
&I(f)\ \leftrightarrow H(f)\land f:A\bij B.
\end{aligned}}
\UebQuellen{\FormulaRefAuto{SemigroupWithZeroHomDef}[def];
\FormulaRefAuto{SemigroupWithZeroIsoDef}[def]}
&
\textbf{Komposition erhält die Struktur}
Für Homomorphismen \(f:A\to B\), \(g:B\to C\) zwischen
Strukturen dieser Klasse ist \(g\circ f\) wieder ein Homomorphismus.
Mit \(\bullet\) als Operation auf \(C\) gilt für \(x,y\in A\):\UebFormel{\begin{aligned}
&(g\circ f)(x\star y)\\
&\quad=(g\circ f)(x)\mathbin{\bullet}(g\circ f)(y).
\end{aligned}}Sind \(f,g\) Isomorphismen, so ist auch \(g\circ f\) ein Isomorphismus.
\UebQuellen{\FormulaRefAuto{SemigroupHomComposition}[thm];
\FormulaRefAuto{SemigroupIsoComposition}[thm]} \\
\addlinespace[.8ex]
\end{BandUebersicht}

